% Part: incompleteness
% Chapter: representability-in-q
% Section: c

\documentclass[../../../include/open-logic-section]{subfiles}

\begin{document}

\olfileid{inc}{req}{cee}
\olsection{دوال~$C$}

لتكن~$C$ أصغر مجموعة من الدوال تحتوي
\begin{enumerate}
\item الصفر،
\item والخلف،
\item والجمع،
\item والضرب،
\item والإسقاطات،
\item والدالة المميزة للمساواة،~$\Char{=}$؛
\end{enumerate}
ومغلقة تحت
\begin{enumerate}
\item التركيب،
\item والبحث غير المحدود المطبق على الدوال المنتظمة.
\end{enumerate}
تذكّر أن هذا القيد الأخير لا يعني إلا أنه لا يجوز استخدام عملية~$\mu$
إلا حين تكون النتيجة كلية. قارن ذلك بتعريف الدوال
\emph{العودية العامة}: فقد أضفنا هنا الجمع والضرب و$\Char{=}$، لكننا
أسقطنا العودية الأولية.

من الواضح أن كل ما في~$C$ عودي، لأن الجمع والضرب و$\Char{=}$ عودية.
وسنبيّن أن العكس صحيح أيضًا؛ وهذا يعني أنه يمكننا تنفيذ العودية
الأولية باستخدام بقية الدوال وعمليات الإغلاق في~$C$.

\end{document}
